\chapter{Valeur extractible maximale}\label{ch:m3:maximal-extractable-value}

Un trader envoie une transaction qui achète pour un million de dollars d'ether dans un pool public, en acceptant jusqu'à
un demi-pour-cent de glissement. Avant son inclusion, un achat d'ether passé par un tiers se place dans le même
bloc, juste devant elle, et une vente de cet ether juste derrière. Le trader reçoit le moins que sa
tolérance autorisait~; l'autre partie empoche l'écart, diminué de ce qu'elle a versé au constructeur du bloc pour
le privilège de ces deux positions. Rien dans la bourse n'a été piraté. La valeur vient du droit
d'ordonner les transactions, droit qui, sur une \omterm{def:m3:blockchains-for-traders:chain}{blockchain}, appartient à celui qui construit le bloc, et qui s'achète et
se vend sur un marché qui lui est propre. Ce chapitre explique ce marché~: d'où provient la valeur extractible, le
sandwich et son arithmétique, les arbitrages et les liquidations qui en sont les cousins bénins, les enchères et
la séparation des rôles par laquelle les blocs sont aujourd'hui construits, et les canaux privés et les dispositifs destinés à
garder les ordres des traders à l'écart de ce marché.

\section{Pouvoir d'ordonnancement et valeur extractible}

\begin{definition}[Valeur extractible maximale, searcher]\label{def:m3:maximal-extractable-value:mev}
\emph{valeur extractible maximale}\index{valeur extractible maximale} (maximal extractable value, MEV) est la valeur qui peut être obtenue en
incluant, en excluant ou en réordonnant des transactions dans un bloc, au-delà des frais ordinaires. Un
\emph{searcher}\index{searcher} est un participant, généralement un système automatisé, qui surveille les transactions
en attente et l'état de la chaîne pour repérer de telles opportunités, et soumet des transactions pour les capter.
\end{definition}

Le terme a été introduit sous le nom de «~miner extractable value~» par Daian et ses coauteurs, qui ont décrit des robots faisant monter
les frais de transaction les uns contre les autres dans des enchères ouvertes, pour être en tête sur les opportunités d'arbitrage sur les
\omterm{def:m3:automated-market-makers:amm}{plateformes décentralisées} (DEX), les ont assimilés à des traders à haute fréquence et ont averti que les frais payés pour la priorité faisaient peser un
risque pour le consensus de la chaîne. La valeur a trois sources. Une transaction peut créer un écart de prix sur lequel
quelqu'un peut opérer (un swap de grande taille déplace un pool). Une variation des prix externes peut en créer un (le prix
périmé d'un \omterm{def:m3:automated-market-makers:amm}{teneur de marché automatisé} (AMM), voir le \cref{ch:m3:automated-market-makers}). Un protocole peut enfin rémunérer celui qui agit le premier (un
bonus de liquidation, voir le \cref{ch:m3:defi-credit-and-leverage}). Dans chaque cas, la valeur revient à celui qui peut passer en premier, ou
à s'intercaler autour de lui, dans l'ordre d'un bloc.

\begin{definition}[Front-running et back-running]\label{def:m3:maximal-extractable-value:front}
\emph{front-running}\index{front-running} consiste à placer une transaction devant une transaction connue en attente, pour profiter de la
variation de prix que cette dernière va provoquer. \emph{back-running}\index{back-running} consiste à placer une transaction
immédiatement après une transaction connue en attente, pour capter l'opportunité qu'elle laisse (un arbitrage, une liquidation) sans
lui nuire.
\end{definition}

\section{Front-running et attaque sandwich}

\begin{definition}[Attaque sandwich, tolérance de glissement]\label{def:m3:maximal-extractable-value:sandwich}
Une \emph{attaque sandwich}\index{attaque sandwich} est un \omterm{def:m3:maximal-extractable-value:front}{front-running} et un \omterm{def:m3:maximal-extractable-value:front}{back-running} autour du swap d'une victime dans un même pool~:
l'attaquant achète ce que la victime achète, celle-ci achète au prix le plus élevé, et l'attaquant revend après elle.
La \emph{tolérance de glissement}\index{tolérance de glissement} de la victime est le pire prix, exprimé en pourcentage sous la
cotation à la soumission, auquel elle autorise l'exécution de son swap~; en deçà, le swap est annulé.
\end{definition}

\begin{omfigure}
\centering
\begin{tikzpicture}[font=\small,
  t/.style={draw,thick,rounded corners=2pt,align=center,minimum height=0.85cm,text width=3.0cm},
  arr/.style={-{Stealth[length=2.2mm]},thick}]
\node[t,draw=omExo,fill=omExo!8] (a) at (0,0) {1.\ l'attaquant achète de l'ETH\\(front-running)};
\node[t,draw=omThm,fill=omThm!8] (v) at (4,0) {2.\ la victime achète de l'ETH\\à son pire prix};
\node[t,draw=omExo,fill=omExo!8] (b) at (8,0) {3.\ l'attaquant vend de l'ETH\\(back-running)};
\draw[arr] (a) -- (v); \draw[arr] (v) -- (b);
\node[font=\footnotesize,anchor=north,align=center,text width=11cm] at (4,-0.8) {un même bloc, un même pool~: le constructeur place les trois
transactions consécutivement~; la tolérance de la victime borne de combien l'étape 1 peut faire monter le prix};
\end{tikzpicture}

\omcaption{Une \omterm{def:m3:maximal-extractable-value:sandwich}{attaque sandwich}. Les deux transactions de l'attaquant encadrent celle de la victime dans le même bloc~; la \omterm{def:m3:maximal-extractable-value:sandwich}{tolérance de glissement} de la victime est
la seule limite à la taille de la première. Schéma de principe.}\label{fig:m3:maximal-extractable-value:sandwich}
\end{omfigure}

\begin{proposition}[La tolérance borne l'attaque sandwich]\label{prop:m3:maximal-extractable-value:bound}
Soit une victime qui échange $\Delta x$ de \omterm{def:m3:blockchains-for-traders:key}{jeton} A contre le \omterm{def:m3:blockchains-for-traders:key}{jeton} B dans un pool à produit constant, avec une sortie minimale égale à $(1 - \tau)$ fois la
cotation à la soumission. La sortie de la victime diminue à mesure que le \omterm{def:m3:maximal-extractable-value:front}{front-running} de l'attaquant croît, de sorte qu'il existe un \omterm{def:m3:maximal-extractable-value:front}{front-running} maximal
$a^*(\tau)$ qui laisse la victime exactement à son minimum~; tout \omterm{def:m3:maximal-extractable-value:front}{front-running} plus grand entraîne l'annulation du swap de la victime
et l'attaque échoue. La perte de la victime est au plus $\tau$ fois sa sortie cotée, et $a^*$ ainsi que le gain brut de l'attaquant croissent avec
$\tau$. Pour $\tau = 0$, aucune \omterm{def:m3:maximal-extractable-value:sandwich}{attaque sandwich} n'est possible.
\end{proposition}

\begin{proof}
Un \omterm{def:m3:maximal-extractable-value:front}{front-running} fait monter le prix de B dans le pool, de sorte que la sortie de la victime (donnée par la \cref{prop:m3:automated-market-makers:swap}) est décroissante
en fonction du \omterm{def:m3:maximal-extractable-value:front}{front-running}~; l'ensemble admissible est un intervalle $[0, a^*]$, avec $a^*$ croissant en $\tau$. La victime reçoit au moins son
minimum, elle perd donc au plus $\tau$ de sa cotation. Pour $\tau = 0$, seul un \omterm{def:m3:maximal-extractable-value:front}{front-running} nul est admissible.
\end{proof}

Dans l'intervalle admissible, le gain de l'attaquant croît avec la taille de son \omterm{def:m3:maximal-extractable-value:front}{front-running}, si bien qu'il attaque à la borne~: le tutoriel
trouve la borne par recherche dichotomique, en arithmétique entière du pool. Pour un achat d'ether de 1~million de dollars dans un pool détenant
15~millions de dollars et 5\,000 ether, avec des frais de 0.30~\%, une tolérance de 0.5~\% autorise un \omterm{def:m3:maximal-extractable-value:front}{front-running} de 38\,912~dollars et un gain brut de
5\,067~dollars~; la victime reçoit 1.56 ether de moins que sa cotation, soit environ 4\,700~dollars au prix initial du pool. À 1~\%, le
\omterm{def:m3:maximal-extractable-value:front}{front-running} double pour atteindre 78\,119~dollars, le gain à 10\,130~dollars et la perte de la victime à 3.12 ether.

\begin{omfigure}
\begin{tikzpicture}
\begin{axis}[width=11cm,height=5cm,
  xlabel={tolérance de glissement de la victime, \%},ylabel={milliers de dollars},
  tick label style={font=\small},label style={font=\small},
  xmin=0,xmax=2,ymin=0,ymax=22,grid=major,grid style={gray!25},
  legend columns=2,legend style={font=\footnotesize,at={(0.5,-0.3)},anchor=north,draw=none,fill=none,/tikz/every even column/.append style={column sep=8pt}},
  table/col sep=comma]
\addplot[omExo,very thick] table[x=tol_pct,y=gross_k]{figdata/markets-3/22-maximal-extractable-value/tolerance.csv};\addlegendentry{gain brut de l'attaquant}
\addplot[omThm,very thick,dashed] table[x=tol_pct,y=loss_k]{figdata/markets-3/22-maximal-extractable-value/tolerance.csv};\addlegendentry{perte de la victime (à 3\,000~dollars par ether)}
\end{axis}
\end{tikzpicture}

\omcaption{L'\omterm{def:m3:maximal-extractable-value:sandwich}{attaque sandwich} d'un achat d'ether de 1~million de dollars dans un pool de 15~millions de dollars et 5\,000 ether (frais de 0.30~\%), au plus grand
\omterm{def:m3:maximal-extractable-value:front}{front-running} que la tolérance de la victime autorise (voir la \cref{prop:m3:maximal-extractable-value:bound}). Les deux grandeurs croissent linéairement avec la tolérance. Arithmétique entière
du pool~; pool illustratif. Données~: le tutoriel du chapitre.}\label{fig:m3:maximal-extractable-value:tolerance}
\end{omfigure}

\begin{omfigure}
\begin{tikzpicture}
\begin{axis}[width=11cm,height=5cm,
  xlabel={front-running, milliers de dollars},ylabel={gain brut de l'attaquant, milliers de dollars},
  tick label style={font=\small},label style={font=\small},
  xmin=0,xmax=78,ymin=0,ymax=10.5,grid=major,grid style={gray!25},unbounded coords=jump,
  legend columns=2,legend style={font=\footnotesize,at={(0.5,-0.3)},anchor=north,draw=none,fill=none,/tikz/every even column/.append style={column sep=8pt}},
  table/col sep=comma]
\addplot[gray,thick,dashed] table[x=front_k,y=gross_k]{figdata/markets-3/22-maximal-extractable-value/curve.csv};\addlegendentry{si la victime n'annulait pas}
\addplot[omExo,very thick] table[x=front_k,y=feasible_k]{figdata/markets-3/22-maximal-extractable-value/curve.csv};\addlegendentry{admissible à une tolérance de 0.5~\%}
\end{axis}
\end{tikzpicture}

\omcaption{Gain brut de l'attaquant en fonction de la taille de son \omterm{def:m3:maximal-extractable-value:front}{front-running} pour le swap décrit à la \cref{fig:m3:maximal-extractable-value:tolerance}
à une tolérance de 0.5~\%. Le gain continue de croître, mais au-delà de 38\,912~USD le swap de la victime serait annulé et l'attaquant
se retrouverait avec de l'ether acheté à prix majoré. Données~: le tutoriel du chapitre.}\label{fig:m3:maximal-extractable-value:curve}
\end{omfigure}

Le calcul indique au trader ce qu'il doit faire. La tolérance est une option vendue à quiconque peut ordonner le bloc~; elle doit être fixée à partir de
l'impact attendu sur le prix de la transaction elle-même, majoré d'une petite marge, et non d'un nombre rond~; les gros swaps doivent être fractionnés ou envoyés par des canaux
où l'ordre n'est pas public (voir ci-dessous).

\section{Back-running, arbitrage CEX-DEX et courses à la liquidation}

La plus grande partie de la valeur extraite ne nuit pas à un utilisateur en particulier.

\begin{definition}[Arbitrage CEX-DEX]\label{def:m3:maximal-extractable-value:cexdex}
\emph{arbitrage CEX-DEX}\index{arbitrage CEX-DEX} consiste à ramener le cours périmé d'un \omterm{def:m3:automated-market-makers:amm}{teneur de marché automatisé} (AMM) au cours des
\omterm{def:m3:centralised-exchanges:cex}{plateformes centralisées} après leur mouvement, le plus souvent couverte aussitôt sur une \omterm{def:m3:centralised-exchanges:cex}{plateforme centralisée} et placée en tête du bloc suivant.
\end{definition}

C'est le \omterm{def:m3:automated-market-makers:lvr}{loss-versus-rebalancing} (LVR) du \cref{ch:m3:automated-market-makers} vu du côté de l'arbitragiste. Dans le \omterm{def:m3:automated-market-makers:amm}{pool de liquidité} du tutoriel, une hausse de 1~\%
du cours de l'ether sur les \omterm{def:m3:centralised-exchanges:cex}{plateformes centralisées} permet à un arbitragiste d'acheter pour environ 52\,000~USD d'ether au pool, avec un profit d'environ 180~USD
après les frais du pool~; une hausse de 2~\%, environ 127\,000~USD avec un profit d'environ 1\,070~USD. Le profit est une course~: celui dont la transaction est la première
dans le bloc l'emporte tout, et la concurrence porte sur le droit d'être le premier. Les courses à la liquidation sont la même course, pour le bonus de liquidation d'un protocole de prêt.
Le \omterm{def:m3:maximal-extractable-value:front}{back-running} d'un gros swap, pour arbitrer le pool qu'il a déplacé, ne nuit pas à celui qui le passe, et les canaux privés ci-dessous
restituent une partie de sa valeur à celui qui passe le swap.

\section{Enchères de gas, séparation proposeur-constructeur et bundles}

\begin{definition}[Enchère de gas prioritaire]\label{def:m3:maximal-extractable-value:pga}
Une \emph{enchère de gas prioritaire}\index{enchère de gas prioritaire} est la compétition ouverte dans laquelle les \omterm{def:m3:maximal-extractable-value:mev}{searchers} surenchérissent sur les \omterm{def:m3:blockchains-for-traders:gas}{frais de priorité} des transactions du
\omterm{def:m3:blockchains-for-traders:mempool}{mempool} public, chacun surenchérissant sur les autres, pour être classé en premier pour une même opportunité.
\end{definition}

Les \omterm{def:m3:maximal-extractable-value:pga}{enchères de gas prioritaires} inondent le réseau de transactions échouées ou remplacées et confèrent le droit d'ordonner au producteur de blocs, qui peut aussi prendre
l'opportunité elle-même. Depuis le passage d'Ethereum à la \omterm{def:m3:blockchains-for-traders:pow}{preuve d'enjeu} (proof of stake), le marché s'est organisé autrement.

\begin{definition}[Séparation proposeur-constructeur, constructeur de blocs, relais MEV, bundle de transactions]\label{def:m3:maximal-extractable-value:pbs}
La \emph{séparation proposeur-constructeur}\index{séparation proposeur-constructeur} est la répartition de la production des blocs entre \emph{constructeurs
de blocs}\index{constructeur de blocs}, des acteurs spécialisés qui assemblent les blocs les plus rentables à partir des transactions publiques et des bundles privés et enchérissent pour
le droit de les faire proposer, et des proposeurs (\omterm{def:m3:blockchains-for-traders:pow}{validateurs}) qui retiennent la meilleure enchère. Un \emph{relais MEV}\index{relais MEV} est un intermédiaire qui
reçoit les blocs des constructeurs, garde leur contenu caché au proposeur tant qu'il ne s'engage pas à proposer, et transmet la meilleure enchère. Un \emph{bundle
de transactions}\index{bundle de transactions} est une liste ordonnée de transactions que soumet un \omterm{def:m3:maximal-extractable-value:mev}{searcher} aux constructeurs, pour être incluse ensemble et dans cet ordre,
ou pas du tout.
\end{definition}

\begin{omfigure}
\centering
\begin{tikzpicture}[font=\small,
  s/.style={draw,thick,rounded corners=2pt,align=center,minimum height=0.9cm,text width=2.2cm},
  arr/.style={-{Stealth[length=2.2mm]},thick}]
\node[s,draw=omDef,fill=omDef!8] (u) at (0,0) {utilisateurs et\\portefeuilles};
\node[s,draw=omExo,fill=omExo!8] (se) at (2.9,0) {searchers\\(bundles de\\transactions)};
\node[s,draw=omThm,fill=omThm!8] (bu) at (5.8,0) {constructeurs\\de blocs};
\node[s,draw=omProp,fill=omProp!8] (re) at (8.7,0) {relais MEV};
\node[s,draw=gray,fill=gray!8] (pr) at (11.6,0) {proposeur\\(validateur)};
\draw[arr] (u) -- (se); \draw[arr] (se) -- (bu); \draw[arr] (bu) -- (re); \draw[arr] (re) -- (pr);
\draw[arr,dashed] (u.south) to[bend right=18] node[below,font=\footnotesize] {flux d'ordres privé} (bu.south);
\end{tikzpicture}

\omcaption{La chaîne d'approvisionnement d'un bloc Ethereum sous \omterm{def:m3:maximal-extractable-value:pbs}{séparation proposeur-constructeur}~: les \omterm{def:m3:maximal-extractable-value:mev}{searchers} envoient des bundles aux constructeurs, qui envoient des blocs et
des enchères aux relais, et le proposeur signe le bloc le plus rentable sans en avoir vu le contenu au préalable. Le \omterm{def:m3:maximal-extractable-value:private}{flux d'ordres privé} va directement aux constructeurs.
Schéma.}\label{fig:m3:maximal-extractable-value:pbs}
\end{omfigure}

\begin{dated}{2026-09}{Séparation proposeur-constructeur en pratique}\label{dat:m3:maximal-extractable-value:boost}
Flashbots décrit MEV-Boost comme un intergiciel open source exploité par des \omterm{def:m3:blockchains-for-traders:pow}{validateurs} pour accéder à un marché concurrentiel de construction de blocs, sa mise en œuvre de
la \omterm{def:m3:maximal-extractable-value:pbs}{séparation proposeur-constructeur} pour l'Ethereum en \omterm{def:m3:blockchains-for-traders:pow}{preuve d'enjeu}~: les constructeurs envoient des blocs aux relais, qui transmettent au proposeur le bloc le plus rentable.
Un témoignage lors d'un procès pénal en 2025 chiffrait à 95~\% la part des \omterm{def:m3:blockchains-for-traders:pow}{validateurs} utilisant MEV-Boost en avril 2023. Les statistiques de Relayscan pour les 24 heures allant jusqu'au 24
septembre 2026 ont compté 6\,526 blocs livrés par des relais, sur environ 7\,200 créneaux, dont trois constructeurs en ont produit environ 93~\%~; le
tableau de bord MEV-Boost tenu par Toni Wahrstätter chiffrait à 91.58~\% la part des blocs MEV-Boost le même jour, avec 7 relais et 25
constructeurs actifs sur quatorze jours.
\end{dated}

Les constructeurs se disputent aux enchères chaque créneau, et les \omterm{def:m3:maximal-extractable-value:mev}{searchers} se disputent la place dans les blocs des constructeurs en les payant~: sur un marché concurrentiel, un
\omterm{def:m3:maximal-extractable-value:mev}{searcher} ne garde qu'une faible part de la valeur d'une opportunité et verse le reste au constructeur, qui en transmet l'essentiel au proposeur. Avec un
versement de 90~\%, l'\omterm{def:m3:maximal-extractable-value:sandwich}{attaque sandwich} de la section précédente rapporte à son \omterm{def:m3:maximal-extractable-value:mev}{searcher} environ 487~USD sur 5\,067~USD, et au constructeur environ 4\,560~USD. La valeur
extraite de la victime revient surtout au \omterm{def:m3:blockchains-for-traders:pow}{validateur}.

Le système repose sur la confiance accordée aux relais et à leur code. Le parquet de New York a poursuivi deux frères pour avoir exploité une faille du logiciel de relais MEV-Boost en
avril 2023~: selon le dossier déposé par l'accusation après le procès, ils ont piégé les bots sandwich d'autres \omterm{def:m3:maximal-extractable-value:mev}{searchers} avec des transactions, enregistré des \omterm{def:m3:blockchains-for-traders:pow}{validateurs} pour
obtenir auprès d'un relais le contenu privé d'un bloc en soumettant un en-tête de bloc invalide, et remplacer les transactions du bundle pour s'emparer de plus de 26
millions de dollars. Le procès a été déclaré nul le 7 novembre 2025, le jury n'ayant pu parvenir à un verdict~; la défense a alors demandé l'acquittement, une requête
plaidée le 23 mars 2026.

\section{Flux d'ordres privé, séquenceurs de rollups et atténuations}

\begin{definition}[Flux d'ordres privé, enchère de flux d'ordres]\label{def:m3:maximal-extractable-value:private}
Le \emph{flux d'ordres privé}\index{flux d'ordres privé} désigne des transactions envoyées directement aux constructeurs ou par des canaux privés plutôt que dans le
\omterm{def:m3:blockchains-for-traders:mempool}{mempool} public, où les \omterm{def:m3:maximal-extractable-value:mev}{searchers} ne peuvent pas les voir. Une \emph{enchère de flux d'ordres}\index{enchère de flux d'ordres} est un mécanisme qui permet aux \omterm{def:m3:maximal-extractable-value:mev}{searchers} d'enchérir pour le droit de
faire un \omterm{def:m3:maximal-extractable-value:front}{back-running} sur la transaction d'un utilisateur, en lui reversant une partie des gains, sans révéler la transaction dans son intégralité.
\end{definition}

\begin{dated}{2026-09}{Un canal privé et une enchère de flux d'ordres}\label{dat:m3:maximal-extractable-value:protect}
Flashbots Protect envoie les transactions vers un \omterm{def:m3:blockchains-for-traders:mempool}{mempool} privé, à l'abri des bots de \omterm{def:m3:maximal-extractable-value:front}{front-running} et sandwich, ne les inclut que si elles n'échouent pas, et
rétrocède la MEV et les \omterm{def:m3:blockchains-for-traders:gas}{frais de priorité}. Son nœud MEV-Share partage avec les \omterm{def:m3:maximal-extractable-value:mev}{searchers} certaines informations sur la transaction d'un utilisateur, simule leurs bundles et
transmet ceux qui aboutissent aux constructeurs à condition que l'utilisateur reçoive une part déterminée de la MEV qu'elles créent (par défaut 90~\%)~; il n'accepte que
les \omterm{def:m3:maximal-extractable-value:front}{back-running}.
\end{dated}

Les canaux privés réduisent les \omterm{def:m3:maximal-extractable-value:sandwich}{attaques sandwich} mais déplacent la confiance~: l'utilisateur compte sur le canal et sur les constructeurs vers lesquels il l'envoie pour ne pas exploiter eux-mêmes le flux,
et un flux qui va vers quelques constructeurs concentre davantage la construction des blocs, ce qui explique en partie que trois constructeurs en produisent la plupart. Sur les \omterm{def:m3:blockchains-for-traders:rollup}{rollups}, le
\omterm{def:m3:blockchains-for-traders:rollup}{séquenceur} (voir le \cref{ch:m3:blockchains-for-traders}) ordonne les transactions~; s'il les traite selon la règle du premier arrivé, premier servi et garde sa file privée, il n'existe
pas de \omterm{def:m3:blockchains-for-traders:mempool}{mempool} public où pratiquer le \omterm{def:m3:maximal-extractable-value:front}{front-running}, et la course se déplace vers la latence jusqu'au \omterm{def:m3:blockchains-for-traders:rollup}{séquenceur}, comme sur une \omterm{def:m3:centralised-exchanges:cex}{plateforme centralisée}. Les conceptions des chapitres précédents traitent le
même problème par l'autre bout~: le \omterm{def:m3:automated-market-makers:routing}{routage par intentions} et les enchères par lots (voir le \cref{ch:m3:automated-market-makers}) exécutent les ordres à un même prix, de sorte que l'ordre
au sein d'un lot n'ait pas d'importance.

\section{Tutoriel~: anatomie d'une attaque sandwich}\label{tut:m3:maximal-extractable-value:sandwich}

\begin{tutorial}
\textbf{Objectif.} Étant donné un pool et le swap d'une victime avec sa tolérance, trouver le \omterm{def:m3:maximal-extractable-value:front}{front-running} le plus grand de l'attaquant, son profit après le \omterm{def:m3:blockchains-for-traders:gas}{gas} et le
paiement du constructeur, et la perte de la victime~; montrer que la tolérance borne les deux. \textbf{État final~:}
les \cref{fig:m3:maximal-extractable-value:tolerance,fig:m3:maximal-extractable-value:curve} et les chiffres du problème du week-end.
\begin{tutsteps}
\item \textbf{L'attaque.} \omterm{def:m3:maximal-extractable-value:front}{Front-running}, victime, \omterm{def:m3:maximal-extractable-value:front}{back-running}, dans l'arithmétique entière du pool~; recherche dichotomique de la borne.
\omcode{code/firm/sandwich/firm_sandwich.py}{26}{61}{L'attaque sandwich au plus grand front-running admissible.}\label{lst:m3:maximal-extractable-value:sandwich}
\item \textbf{Tolérances.} \texttt{by\_tolerance()} avec et sans un versement de 90~\% au constructeur.
\item \textbf{Détection.} Repérer le motif dans les swaps d'un bloc~: un même expéditeur encadre la transaction d'un autre dans le même pool.
\omcode{code/firm/sandwich/firm_sandwich.py}{64}{77}{Détection d'une attaque sandwich dans la liste des swaps d'un bloc.}\label{lst:m3:maximal-extractable-value:detect}
\item \textbf{Exécuter} \texttt{cex\_dex\_arbitrage(3030)} et \texttt{fig\_mev.py}.
\end{tutsteps}
\textbf{Pour aller plus loin.} Laisser la victime fractionner son swap en dix sur plusieurs blocs~; ajouter un \omterm{def:m3:maximal-extractable-value:mev}{searcher} concurrent qui enchérit sur la même \omterm{def:m3:maximal-extractable-value:sandwich}{attaque sandwich} et trouver la
part du constructeur à partir de laquelle l'attaque cesse d'être rentable~; exécuter \texttt{find\_sandwiches} sur un bloc synthétique comportant deux attaques qui se chevauchent.
\end{tutorial}

\section{Construction~: le simulateur de sandwich}\label{bld:m3:maximal-extractable-value:sandwich}

\begin{build}
\bfield{Rôle} La mini-entreprise échange sur la chaîne en tant que client et réalise des arbitrages en tant que \omterm{def:m3:maximal-extractable-value:mev}{searcher}~: elle doit savoir ce que ses propres transactions exposent,
fixer des tolérances, valoriser les back-runs et détecter quand elle a été victime d'une \omterm{def:m3:maximal-extractable-value:sandwich}{attaque sandwich}.
\bfield{Interface} \texttt{run(ra, rb, dx\_v, dx\_a)}~; \texttt{min\_out(ra, rb, dx\_v, tolerance\_bps)}~; \texttt{max\_front\_run}~; \texttt{sandwich(ra,
rb, dx\_v, tolerance\_bps, gas\_cost, builder\_share\_bps)} renvoyant le front-run, la sortie et la perte de la victime, ainsi que les gains bruts et nets~;
\texttt{find\_sandwiches(txs)}. Swaps issus de \texttt{firm.amm}.
\bfield{Règles} Unités de base entières et arrondi exact du pool~; minimum de la victime calculé à partir de la cotation au moment de la soumission, comme le font les portefeuilles~; une attaque
est refusée si la transaction de la victime serait annulée.
\bfield{Tests d'acceptation} \texttt{code/firm/sandwich/tests/}~: la borne est exacte à une unité près~; le gain et la perte croissent avec la tolérance~; la part du constructeur réduit
le net~; une tolérance nulle ne laisse rien~; un sandwich détecté dans un bloc et non entre plusieurs pools.
\bfield{Pour aller plus loin} Pools à \omterm{def:m3:automated-market-makers:concentrated}{liquidité concentrée} franchissant des ticks~; sandwiches multi-pools~; valorisation des back-runs pour une \omterm{def:m3:maximal-extractable-value:private}{enchère de flux d'ordres}.
\end{build}

\begin{omsources}
\item P.~Daian et al., «~Flash Boys 2.0: Frontrunning, Transaction Reordering, and Consensus Instability in Decentralized Exchanges~»,
arXiv:1904.05234, 2019.
\item Flashbots documentation (MEV-Boost, MEV-Share, Flashbots Protect), dépôt GitHub flashbots/flashbots-docs, septembre 2026~; relayscan.io
et mevboost.pics, consulté le 24 septembre 2026.
\item United States v.\ Peraire-Bueno, No.\ 1:24-cr-00293 (S.D.N.Y.)~: mémorandum du gouvernement du 16 janvier 2026 (Doc 236) et la déclaration
de mistrial du 7 novembre 2025 figurant au dossier~; requête conjointe aux fins d'acquittement (Doc 225, 12 décembre 2025)~; ordonnance fixant l'audience de plaidoirie (Doc 240).
\end{omsources}

\section{Exercices}

\begin{exercise}[$\star$]\label{exo:m3:maximal-extractable-value:1}
Expliquer pourquoi personne ne peut lancer une \omterm{def:m3:maximal-extractable-value:sandwich}{attaque sandwich} sur un swap envoyé avec une \omterm{def:m3:maximal-extractable-value:sandwich}{tolérance de glissement} nulle, et pourquoi les portefeuilles ne prennent pas zéro par défaut.
\end{exercise}

\begin{exercise}[$\star$]\label{exo:m3:maximal-extractable-value:2}
Classer comme \omterm{def:m3:maximal-extractable-value:front}{front-running}, \omterm{def:m3:maximal-extractable-value:front}{back-running} ou aucun des deux~: une liquidation d'un prêt sous-collatéralisé placée juste après une mise à jour de l'\omterm{def:m3:blockchains-for-traders:contract}{oracle}~; un arbitrage
placé juste après un swap de grande taille~; un achat placé juste avant un achat de grande taille.
\end{exercise}

\begin{exercise}[$\star$]\label{exo:m3:maximal-extractable-value:3}
Expliquer ce qu'un \omterm{def:m3:maximal-extractable-value:pbs}{relais MEV} garantit au constructeur, et ce qu'il garantit au proposeur.
\end{exercise}

\begin{exercise}[$\star\star$]\label{exo:m3:maximal-extractable-value:4}
La tolérance de la victime double, passant de 0.5~\% à 1~\%. Calculer de combien varient le gain brut de l'attaquant et la perte de la victime.
\end{exercise}

\begin{exercise}[$\star\star$]\label{exo:m3:maximal-extractable-value:5}
Un \omterm{def:m3:maximal-extractable-value:mev}{searcher} verse au constructeur 90~\% du gain brut et 20~dollars de \omterm{def:m3:blockchains-for-traders:gas}{gas}. Calculer le gain net sur l'\omterm{def:m3:maximal-extractable-value:sandwich}{attaque sandwich} à 1~\%.
\end{exercise}

\begin{exercise}[$\star\star$]\label{exo:m3:maximal-extractable-value:6}
Expliquer pourquoi le \omterm{def:m3:maximal-extractable-value:private}{flux d'ordres privé} tend à concentrer la construction des blocs.
\end{exercise}

\begin{exercise}[$\star\star\star$]\label{exo:m3:maximal-extractable-value:7}
\emph{Programmation.} Avec \texttt{cex\_dex\_arbitrage}, déterminer la plus petite hausse du prix sur la \omterm{def:m3:centralised-exchanges:cex}{plateforme centralisée} qui permet à l'arbitrage de couvrir 20~dollars de \omterm{def:m3:blockchains-for-traders:gas}{gas}.
\end{exercise}

\begin{exercise}[$\star\star\star$]\label{exo:m3:maximal-extractable-value:8}
\emph{Repérer l'erreur.} «~Nous envoyons nos swaps par un canal privé, donc nous ne payons aucune MEV.~»
\end{exercise}

\section{Problème~: anatomie d'un sandwich}

\begin{problem}[{Problème du week-end --- un swap de 1~million de dollars et le searcher qui l'a vu}]\label{pb:m3:maximal-extractable-value:1}
Un trader achète de l'ether avec 1~million de dollars dans un pool à produit constant contenant 15~millions de dollars et 5\,000 ether, avec des frais de 0.30~\%, via le
\omterm{def:m3:blockchains-for-traders:mempool}{mempool} public. Les deux transactions du \omterm{def:m3:maximal-extractable-value:mev}{searcher} coûtent 20~dollars de \omterm{def:m3:blockchains-for-traders:gas}{gas}.

\medskip\noindent\textbf{Partie~I --- La cotation.}
\begin{enumerate}
\item Combien d'ether la cotation promet-elle, et quel est l'impact propre du swap sur le prix~?
\item Quelle est la sortie minimale du trader à une tolérance de 0.5~\%~? Et à 1~\%~?
\item Quel est le front-run le plus grand à chaque tolérance~?
\item Pourquoi est-ce le meilleur choix du \omterm{def:m3:maximal-extractable-value:mev}{searcher}~?
\item Que ferait le \omterm{def:m3:maximal-extractable-value:mev}{searcher} avec une tolérance de 0~\%~?
\end{enumerate}

\medskip\noindent\textbf{Partie~II --- Le profit.}
\begin{enumerate}[resume]
\item Quel est le gain brut du \omterm{def:m3:maximal-extractable-value:mev}{searcher} à chaque tolérance~?
\item Combien le trader perd-il, en ether et à 3\,000 dollars par ether~?
\item Quel gain net le \omterm{def:m3:maximal-extractable-value:mev}{searcher} obtient-il s'il ne verse rien au constructeur~? Et s'il en verse 90~\%~?
\item Qui finit par capter la plus grande part de la valeur~?
\item Pourquoi un marché concurrentiel des constructeurs pousse-t-il la part du \omterm{def:m3:maximal-extractable-value:mev}{searcher} vers zéro~?
\end{enumerate}

\medskip\noindent\textbf{Partie~III --- Défenses.}
\begin{enumerate}[resume]
\item Quelle tolérance faut-il fixer pour ce swap, et pourquoi~?
\item Que changerait le fractionnement du swap en dix tranches réparties sur dix blocs~?
\item Que changerait un canal privé, et que ne changerait-il pas~?
\item Comment un \omterm{def:m3:automated-market-makers:routing}{routage par intentions} ou une enchère par lots modifierait-il le résultat~?
\item Comment détecter après coup qu'un swap a été pris en sandwich~?
\end{enumerate}

\medskip\noindent\textbf{Partie~IV --- Jugement.}
\begin{enumerate}[resume]
\item Une \omterm{def:m3:maximal-extractable-value:sandwich}{attaque sandwich} constitue-t-elle un \omterm{def:m3:maximal-extractable-value:front}{front-running} au sens des règles d'abus de marché qui s'appliquent aux courtiers sur les marchés réglementés~?
\item Pourquoi les back-runs et l'\omterm{def:m3:maximal-extractable-value:cexdex}{arbitrage CEX-DEX} sont-ils traités différemment des \omterm{def:m3:maximal-extractable-value:sandwich}{attaques sandwich}~?
\item Que révèle l'exploitation du relais de 2023 sur la chaîne d'approvisionnement de la construction des blocs~?
\item Énoncer le \emph{résultat nommé}~: le profit net maximal du \omterm{def:m3:maximal-extractable-value:mev}{searcher} et la perte du trader à des tolérances de 0.5~\% et de 1~\%.
\item En une phrase~: qu'est-ce qu'une \omterm{def:m3:maximal-extractable-value:sandwich}{tolérance de glissement}~?
\end{enumerate}
\end{problem}

\section{Questions d'entretien}

\begin{interviewq}[$\star$ \iqroles{trader}]\label{iq:m3:maximal-extractable-value:1}
Qu'est-ce qu'une \omterm{def:m3:maximal-extractable-value:sandwich}{attaque sandwich}, et comment protégez-vous vos swaps contre ce type d'attaque~?
\end{interviewq}

\begin{interviewq}[$\star$ \iqroles{développeur}]\label{iq:m3:maximal-extractable-value:2}
Expliquez la \omterm{def:m3:maximal-extractable-value:pbs}{séparation proposeur-constructeur} (PBS) et le rôle des \omterm{def:m3:maximal-extractable-value:pbs}{relais MEV}.
\end{interviewq}

\begin{interviewq}[$\star\star$ \iqroles{chercheur}]\label{iq:m3:maximal-extractable-value:3}
Dérivez le front-run le plus grand que la tolérance de la victime autorise dans un pool à produit constant.
\end{interviewq}

\begin{interviewq}[$\star\star$ \iqroles{trader, chercheur}]\label{iq:m3:maximal-extractable-value:4}
Quel lien existe-t-il entre l'\omterm{def:m3:maximal-extractable-value:cexdex}{arbitrage CEX-DEX} et le \omterm{def:m3:automated-market-makers:lvr}{loss-versus-rebalancing} (LVR) d'un pool~?
\end{interviewq}

\begin{interviewq}[$\star\star$ \iqroles{risque}]\label{iq:m3:maximal-extractable-value:5}
Quels risques une entreprise qui exploite des \omterm{def:m3:maximal-extractable-value:mev}{searchers} prend-elle, sur les plans opérationnel et juridique~?
\end{interviewq}

\begin{interviewq}[$\star\star\star$ \iqroles{développeur, chercheur}]\label{iq:m3:maximal-extractable-value:6}
Concevez une \omterm{def:m3:maximal-extractable-value:private}{enchère de flux d'ordres} qui restitue de la valeur aux utilisateurs sans divulguer leurs transactions.
\end{interviewq}
